% Excerpt from K. F. Pal and T. Vertesi, 'Maximal violation of a bipartite three-setting, two-outcome
% Bell inequality using infinite-dimensional quantum systems', Phys. Rev. A 82, 022116 (2010),
% arXiv:1006.3032v1 (source I3322inf.tex), Section 'Other examples': Table I and the text discussing it.
% Supersedes the 2009 windows (found by the auditor agent on 2026-10-11 23:49).

\begin{table*}[t]
\caption{Results for the $19$ Bell inequalities for which the tightness of
the quantum bound has not been proven. The local Hilbert space necessary to achieve the best
lower bound is denoted by $\cal H$, column 'New' tells if this lower bound is an improvement
over our previous result, and column 'Level' is the level to derive the upper bound.}
\vskip 0.2truecm
\centering
\begin{tabular}{l c l l c l}
\hline\hline

Case\hphantom{W}&Lower bound&$\cal H$&New&Upper bound-&\hphantom{ww}Level\\
&&&&lower bound&\\
 \hline
$I^{20}_{4422}$&0.4676794&$\RR^{4}$&no&0.0000000&$L3+aa'bb'$\\
$J^{30}_{4422}$&0.4363842&$\RR^{17}$&yes&0.0008787&$L3$\\
$J^{34}_{4422}$&0.4492657&$\RR^{4}$&no&0.0000000&$L3+aa'bb'$\\
$J^{48}_{4422}$&0.7516220&$\RR^{12}$&yes&0.0000120&$L3+aa'bb'$\\
$J^{62}_{4422}$&0.7500755&$\RR^{6}$&no&0.0000166&$L3$\\
$J^{73}_{4422}$&1.1223170&$\RR^{20}$&yes&0.0002409&$L3$\\
$J^{97}_{4422}$&1.1583626&$\RR^{2}$&no&0.0001425&$L3$\\
$A_{14}$&0.4759513&$\RR^{20}$&yes&0.0018953&$L3$\\
$A_{21}$&0.3260601&$\RR^{10}$&yes&0.0001176&$L3$\\
$A_{47}$&0.4608544&$\CC^{2}$&no&0.0020847&$L2+aa'b+abb'$\\
$A_{62}$&0.4065268&$\RR^{15}$&yes&0.0003457&$L2+aa'b+abb'$\\
$A_{64}$&0.3900890&$\RR^{3}$&no&0.0013878&$L2+aa'b+abb'$\\
$A_{65}$&0.3688996&$\RR^{5}$&yes&0.0000000&$L2+aa'b+abb$\\
$A_{67}$&0.3990671&$\RR^{5}$&no&0.0000000&$L3$\\
$A_{68}$&0.4025522&$\RR^{18}$&yes&0.0024807&$L2+aa'b+abb'$\\
$A_{80}$&0.3769863&$\RR^{4}$&no&0.0002557&$L2+aa'b+abb'$\\
$A_{82}$&0.4708838&$\RR^{20}$&yes&0.0000341&$L2+aa'b+abb'$\\
$A_{84}$&0.6352087&$\RR^{18}$&yes&0.0000747&$L2+aa'b+abb'$\\
$A_{89}$&0.3035637&$\RR^{17}$&yes&0.0022575&$L2+aa'b+abb'$\\
 \hline
 \end{tabular}
 \label{table:highlowdiff}
 \end{table*}

In Ref.~\cite{PV09} we calculated both lower and upper bounds for all known tight binary bipartite Bell
inequalities with up to five measurement settings per party. For most cases we managed to
find explicit solutions saturating the upper bound. There were 20 exceptions. The case of
$I_{3322}$ was discussed in the previous chapter. Now we have applied the iterative algorithm
with up to 25 dimensional component spaces for the remaining 19 problems.
In 11 cases we have found better solution than the ones we reported in Ref.~\cite{PV09}. The solution for $A_{65}$ could have been found by our previous methods, but we missed it. All the other ones involve larger Hilbert spaces than those methods could handle. We have also
made the upper bound tighter by going to higher levels by using a computer with larger
memory. As previously, the calculations have been done with Borcher's code CSDP for
semidefinite programming~\cite{CSDP}.
For cases with five measurement settings per party we could do the calculation on level
two plus $aa'b$ plus $abb'$. For the smaller cases we could afford to go up to level
three, except for $I_{3322}$, where we did level four. The notion of levels and partial levels,
and their notation is explained in Ref.~\cite{NPA2}, and also in Ref~\cite{PV09}. This way the two bounds met for $A_{65}$. We sent $A_{67}$ at level three, and $J^{48}_{4422}$ and $J^{97}_{4422}$ at
level three plus $aa'bb'$ to Brian Borchers, as test cases for his new code which requires
much less memory than CSDP for very large cases. For $A_{67}$ this new upper bound does
agree with the lower one. Unfortunately, for $J^{48}_{4422}$ there remained a difference of
$0.000012$, significantly less than the level three value of $0.0000657$, but more than numerical
uncertainty. At an even higher level it would probably disappear, but it is not sure. Therefore,
we can not state that our solution with 12-dimensional component spaces is already the optimum one.
The case $J^{97}_{4422}$ at level three plus $aa'bb'$ seems to be a difficult case for the new
code, it could solve it only at a reduced accuracy, and the value it has given for the
upper bound turned out to be lower than the lower bound. Therefore, very probably the
solution with a pair of real qubits we found the best so far is actually the optimum.
If we added to level three just one quarter of the $aa'bb'$ type terms, we could use CSDP.
In the cases of $I^{20}_{4422}$ and $J^{34}_{4422}$ this was already enough to reach
the lower bound. By looking at the dependence of the upper bound on the level it has been
calculated, we are almost sure that for $J^{30}_{4422}$, $J^{62}_{4422}$, $A_{47}$,
$A_{62}$, $A_{64}$, $A_{82}$ and $A_{89}$ the present lower bound will not be reached,
therefore we do not yet have the optimum solution. Probably $I_{3322}$  is not the only
inequality requiring infinite dimensional Hilbert spaces for maximum violation. Actually, the
difference in the cases of $J^{62}_{4422}$ and $A_{82}$ is very small, but it does not
change much with increasing level. For $A_{80}$ and $A_{84}$ the upper bound would more
probably converge to our solution, while for the rest of the cases we can not tell from the
present results.
